\section{Preference Corpus Curation}
\label{sec:data-construction}

This section details the curation of our preference corpus. We begin by formalizing the task to clarify the data requirements, followed by the collection and augmentation processes.


\subsection{Task Definition}
\label{sec:task_definition}

We model the data-centric task as a pairwise selection task: given a task description, a data report, and two candidate solutions, the objective is to identify the superior solution and estimate a confidence score (Figure~\ref{fig:main}(a)).
Formally, the input $\mathcal{X}$ is:
\begin{equation}
    \mathcal{X} = \left( I, D_{rep}, \{C_0, C_1\}, \mathcal{P} \right)
\end{equation}
where $I$, $D_{rep}$, $\{C_0, C_1\}$, and $\mathcal{P}$ denote the task, data report, code pair, and system prompt, respectively. The output $\mathcal{Y}$ is defined as:
\begin{equation}
    \mathcal{Y} = \left\{ (cot, \hat{y}, c) \mid cot, \; \hat{y} \in \{0, 1\}, \; c \in [0, 1.0] \right\}
\end{equation}
consisting of the reasoning $cot$, predicted winner $\hat{y}$, and confidence $c$, which serves as the gating threshold in Section~\ref{sec:agent}.


\subsection{Source and Scope}
\label{sec:source_scope}

To instantiate the task inputs defined above, we construct a large-scale corpus derived from the real-world execution trajectories of two ML agents, \textbf{AIDE}~\cite{jiang2025aideaidrivenexplorationspace} and \textbf{AutoMind}~\cite{ou2025automindadaptiveknowledgeableagent}, operating on \textbf{MLE-bench}~\cite{chan2025mlebenchevaluatingmachinelearning} platform (Figure~\ref{fig:main}(b)).
Powered by DeepSeek-V3.1~\cite{deepseekai2025deepseekv3technicalreport} and o3-mini~\cite{openai2025o3card}, these agents generate \textbf{1,329} valid solutions across \textbf{26} diverse tasks (Table~\ref{tab:task_distribution}).
Unlike synthetic snippets, these candidates represent \textit{complete ML workflows} that are entirely generated by agents and absent from any pre-training corpora, heavily incorporating logically incomplete but executable intermediate states to model noisy real-world exploration (see Appendix \ref{sec:trajectory_sampling}).
Therefore, identifying the superior solution requires evaluating \textit{how well an algorithm fits the specific data characteristics}, rather than merely checking for code syntax.

\subsection{Dataset Curation and Instantiation}
\label{sec:data_construction}

For rigorous evaluation, we implement an \textbf{Expert-in-the-Loop} pipeline to prune raw trajectories into \textbf{895} high-quality instances. 
This process involves deduplication, taxonomy tagging, and expert sampling to cap dominant methods and ensure algorithmic diversity. Next, we instantiate the dataset by exhaustively generating pairwise combinations from this curated corpus. 
We apply strict filtering to discard ambiguous pairs and balance the ground-truth winner's position to mitigate position bias~\cite{Shi2024judingthejudges}. 
This yields a final dataset of \textbf{18,438 comparisons} (Figure~\ref{fig:main}(c)), utilizing \textbf{micro-averaged accuracy} as the primary metric.


\subsection{Input Augmentation: The Verified Data Analysis Report}
\label{sec:input_augmentation}

To address LLMs' numerical limitations~\cite{davies2025languagemodelsembednumbers, li-etal-2025-exposing} and context constraints preventing direct data ingestion, we augment inputs with a \textbf{Verified Data Analysis Report} that transforms raw statistics into semantic narratives~\cite{Rytting2021leveragingtheinductivebias, zhang2025srllmrethinkingstructuredrepresentation}.
To guarantee factual grounding and prevent hallucination, we implement a strict protocol (Figure~\ref{fig:main}(d)) consisting of three concrete steps:
(1) \textbf{Code Generation}: GPT-5.1~\cite{openai2025gpt5card} generates a Python script to profile raw data with labels masked.
For example, it writes ``\texttt{print(df['target'].value\_counts())}'' to inspect the target distribution.
(2) \textbf{Execution and Verification}: The script runs in a sandbox to produce standard output.
A human expert performs a strict pass or fail validity check to ensure the log is free of runtime errors.
For instance, the execution log returns a raw fact like ``\texttt{Target Distribution: 0: 0.915, 1: 0.085}''.
(3) \textbf{Verbalization}: GPT-5.1 reads this execution log and translates it into a semantic insight.
Following the previous example, it produces the final report: ``Data Imbalance Warning: Severe class imbalance (Pos: 8.5\%). Implication: Accuracy is not a suitable metric; consider using F1-score.''
This process ensures reliable semantic grounding for the task (see case in Appendix Figure~\ref{fig:report_patent}).